\section{Technical Architecture}
\label{sec:architecture}

The architecture of NavHAL is designed to address the "Abstraction-Efficiency Paradox" in embedded systems: the need for high-level, portable APIs that does not introduce the overhead typically associated with heavyweight middleware. This is achieved through a meticulously layered design that decouples peripheral logic from hardware topology.

\subsection{The Layered Abstraction Model}
NavHAL is designed as a vertical abstraction stack that ensures hardware independence while maintaining minimal execution overhead. The architecture is organized into four logical tiers:

\begin{figure}[h]
    \centering
    \begin{tikzpicture}[
        layer/.style={rectangle, draw=navblue, fill=navblue!5, thick, minimum width=8cm, minimum height=0.8cm, text centered, font=\sffamily\bfseries},
        arrow/.style={-stealth, thick, navblue},
        node distance=0.5cm
    ]
        % Nodes
        \node (app) [layer] {1. Application Layer (main.c)};
        \node (common) [layer, below=of app] {2. Common API Layer (include/common)};
        \node (arch) [layer, below=of common] {3. Processor Implementation Layer (src/core)};
        \node (ral) [layer, below=of arch] {4. Register Abstraction Layer (RAL)};
        \node (hw) [layer, fill=gray!20, draw=gray, below=of ral] {Target Hardware (Silicon)};

        % Arrows
        \draw [arrow] (app) -- (common);
        \draw [arrow] (common) -- (arch);
        \draw [arrow] (arch) -- (ral);
        \draw [arrow] (ral) -- (hw);

        % Annotations
        \node [right=0.5cm of app, font=\small\itshape, text=navblue] {Architecture Macro};
        \node [right=0.5cm of common, font=\small\itshape, text=navblue] {Unified Interface};
        \node [right=0.5cm of arch, font=\small\itshape, text=navblue] {Architecture Drivers};
        \node [right=0.5cm of ral, font=\small\itshape, text=navblue] {Struct-based Maps};

    \end{tikzpicture}
    \caption{The NavHAL Abstraction Stack}
    \label{fig:abstraction_stack}
\end{figure}

\begin{enumerate}
    \item \textbf{Application Layer}: The entry point where the target architecture is defined (e.g., \texttt{\#define CORTEX\_M4}) and generic HAL functions are invoked.
    \item \textbf{Common API Layer}: Architecture-agnostic headers in \texttt{include/common/} that route calls to the specific architecture implementation.
    \item \textbf{Processor Implementation Layer}: Source files in \texttt{src/core/<processor>/} that provide the actual logic for peripherals like GPIO, UART, and I2C.
    \item \textbf{Register Abstraction Layer (RAL)}: Low-level headers defining memory-mapped register structures for precise hardware control.
\end{enumerate}

\subsection{The Architecture-Led Dispatch Mechanism}
NavHAL's portability is driven by an explicit architecture selection mechanism. Instead of complex runtime auto-detection, the HAL relies on a clear, preprocessor-driven dispatch system.

\subsubsection{Explicit Architecture Selection}
Selection of the hardware target is performed at the top of the application or via build-system compiler flags. By defining a macro such as \texttt{CORTEX\_M4}, the system "activates" the corresponding low-level headers.

\subsubsection{Header Dispatch Logic}
The headers in the \texttt{include/common/} directory act as the primary dispatchers. They check for defined architecture macros and include the relevant sub-headers from the \texttt{include/core/} directory.

\begin{figure}[h]
    \centering
    \begin{tikzpicture}[
        node/.style={rectangle, draw=navblue, fill=navblue!5, thick, text centered, font=\sffamily\small},
        decision/.style={diamond, draw=navblue, fill=orange!10, thick, text centered, inner sep=0pt, font=\sffamily\tiny, width=2.5cm},
        arrow/.style={-stealth, thick, navblue}
    ]
        % Nodes
        \node (app) [node] {Application (\texttt{main.c})};
        \node (def) [node, below=0.8cm of app, fill=green!5] {\texttt{\#define CORTEX\_M4}};
        \node (common) [node, below=0.8cm of def] {\texttt{include/common/hal\_gpio.h}};
        \node (dispatch) [decision, below=1.0cm of common] {\texttt{if CORTEX\_M4}};
        \node (impl) [node, right=1.5cm of dispatch] {\texttt{include/core/cortex-m4/gpio.h}};

        % Arrows
        \draw [arrow] (app) -- (def);
        \draw [arrow] (def) -- (common);
        \draw [arrow] (common) -- (dispatch);
        \draw [arrow] (dispatch) -- node[above, font=\tiny] {Yes} (impl);

    \end{tikzpicture}
    \caption{Preprocessing Dispatch Flow}
    \label{fig:dispatch_flow}
\end{figure}

\subsubsection{Source File Selection via CMake}
While the headers are dispatched via the preprocessor, the source files are selected by the CMake build system. The root \texttt{CMakeLists.txt} uses the \texttt{CMAKE\_SYSTEM\_PROCESSOR} variable to include the appropriate implementation directory within \texttt{src/core/}. 

This dual-mechanism ensures that only the necessary code is included in the project, fulfilling the \textit{Zero-Cost Abstraction} principle: the hardware-specific implementation is effectively "linked" directly to the common API calls without intermediate layers or runtime overhead.

\subsection{Directory Structure and Organization}
The project follows a strict hierarchical layout to ensure scalability:
\begin{itemize}
    \item \textbf{\texttt{include/common/}}: Architecture-independent API headers.
    \item \textbf{\texttt{include/core/<processor>/}}: Hardware-specific register maps and function prototypes.
    \item \textbf{\texttt{src/core/<processor>/}}: Implementation logic grouped by peripheral (GPIO, UART, etc.).
    \item \textbf{\texttt{src/board/<board>/}}: Board-specific files such as linker scripts and system configuration.
\end{itemize}

\subsection{Registry Abstraction: The Structure-Based Approach}
Unlike many HALs that rely on a sea of \texttt{\#define} constants for register offsets, NavHAL employs a structure-based mapping system. Each peripheral is represented as a packed C \texttt{struct}, where the members correspond to specific hardware registers at precise memory offsets.

\begin{lstlisting}[language=C, caption=GPIO Register Mapping (Cortex-M4)]
typedef struct {
  __IO uint32_t MODER;   /* Mode register */
  __IO uint32_t OTYPER;  /* Output type register */
  __IO uint32_t OSPEEDR; /* Output speed register */
  __IO uint32_t PUPDR;   /* Pull-up/pull-down register */
  /* ... Additional registers ... */
} GPIOx_Typedef;

#define GPIOA_BASE_ADDR 0x40020000
#define GPIOA ((GPIOx_Typedef *)GPIOA_BASE_ADDR)
\end{lstlisting}

This methodology provides several advantages:
\begin{itemize}
    \item \textbf{Type Safety}: Prevents accidental writes to incorrect memory locations by providing structured access.
    \item \textbf{Readability}: Allows code like \texttt{GPIOA->MODER = 0x1;} which is far more intuitive than complex bitwise offset arithmetic.
    \item \textbf{Debuggability}: Modern debuggers (GDB/Ozone) can naturally parse these structures, allowing for real-time inspection of peripheral states in a human-readable format.
\end{itemize}

\subsection{Memory Topology and Port Mapping}
To maintain a unified API for heterogeneous hardware, NavHAL uses an "Absolute Pin Mapping" system. A single integer or enum defines a pin (e.g., \texttt{PA5}), and the HAL provides macros to compute the corresponding port structure and pin offset at compile or runtime.

\begin{lstlisting}[language=C, caption=Port Resolution Macro]
#define GPIO_GET_PORT(pin) ((GPIOx_Typedef *)(GPIOA_BASE_ADDR + ((pin / 16) * 0x400)))
#define GPIO_GET_PIN_OFFSET(pin) (pin % 16)
\end{lstlisting}

This abstraction allows for high-level functions like \texttt{hal\_gpio\_digitalwrite(PA5, HIGH)} to remain consistent across different chip families, even when the underlying port counts or memory layouts differ.

\subsection{Modular Core Separation}
The project directory structure mirrors this architectural philosophy, ensuring that adding support for a new architecture does not require modification of existing core logic:

\begin{itemize}
    \item \texttt{include/common/}: Universal interfaces.
    \item \texttt{include/core/<arch>/}: Architecture-specific definitions.
    \item \texttt{src/core/<arch>/<peripheral>/}: Implementation of the HAL for that architecture.
\end{itemize}

By enforcing this strict modularity, NavHAL prevents "cross-pollution" between targets, keeping the codebase clean and maintainable as the number of supported MCUs grows.
